\uuid{g70E}
\titre{Coefficient}
\auteur{bodin ; hanani ; mzari}
\organisation{exo7}
\licence{CC BY-NC-SA 4.0}
\datecreate{2020-01-20}
\niveau{L1}
\module{Algèbre}
\chapitre{Polynôme, fraction rationnelle}
\sousChapitre{Définition, degré, produit}
\difficulte{3}
\theme{Polynômes, coefficient}
\exo7qcmid{lille-1.polynomes.006}

\contenu{
\begin{question}
\qclassification{105.05}
\qtags{motcle=coefficient}
Soit $P(X) = X^n + a_{n-1}X^{n-1} + \cdots + a_1X+a_0$ un polynôme de $\Rr[X]$ de degré $n \ge 1$. À ce polynôme $P$ on associe un nouveau polynôme $Q$, défini par $Q(X) = P(X - \frac{a_{n-1}}{n})$.

Quelles sont les assertions vraies ?
\begin{answers}
    \bad{Si $P(X) = X^2+3X+1$ alors $Q(X) = X^2-2X$.}

    \good{Si $P(X) = X^3-3X^2+2$ alors $Q(X) = X^3-3X$.}

    \bad{Le coefficient constant du polynôme $Q$ est toujours nul.}

    \good{Le coefficient du monôme $X^{n-1}$ de $Q$ est toujours nul.}
\end{answers}
\begin{explanations}
Cette transformation est faite afin que le coefficient du monôme $X^{n-1}$ de $Q$ soit toujours nul.
\end{explanations}
\end{question}
}
